\documentclass[12pt,a4paper]{article}
\usepackage{fontspec}
% Prefer installed fonts; bundled unmodified fonts keep CI and local builds portable.
\IfFontExistsTF{Liberation Sans}{\setmainfont{Liberation Sans}}{%
  \setmainfont{LiberationSans}[Path=fonts/,Extension=.ttf,UprightFont=*-Regular,BoldFont=*-Bold,ItalicFont=*-Italic,BoldItalicFont=*-BoldItalic]}

\usepackage{babel}
\babelprovide[main,import]{portuguese}
\usepackage{hyperref,booktabs,array,xcolor}
\hypersetup{colorlinks=true,urlcolor=blue,linkcolor=black}
\usepackage[margin=2.5cm]{geometry}
\title{EstudaOffline: implementação e verificação técnica de um MVP local em Flutter web}
\author{Estudantes — identificação omitida nesta demonstração}
\date{Outubro de 2026}
\begin{document}\sloppy\maketitle
\begin{abstract}Este trabalho descreve um MVP de planejamento de estudos com dados locais e um protocolo de verificação de persistência, backup e carregamento offline. O estudo usa dados sintéticos e cenários técnicos, sem participantes humanos. Os resultados são apresentados somente quando sustentados pelos registros de execução. A contribuição é um artefato e um procedimento verificável; não se demonstra eficácia pedagógica.\end{abstract}
\noindent\textbf{Palavras-chave:} Flutter web; PWA; armazenamento local; testes; planejamento de estudos.
\section{Introdução}O EstudaOffline é um produto mínimo viável (MVP) de planejamento de estudos em Flutter para a web, com armazenamento local e proposta de funcionamento como aplicação web progressiva (PWA). O problema técnico é manter tarefas acessíveis no navegador quando a conexão está indisponível, reconhecendo limites de cache, armazenamento e atualização. O contexto escolar é o CETI Lucas Meireles Alves, curso técnico em Análise e Desenvolvimento de Sistemas, Teresina, bairro Chapadinha Sul. O objetivo é implementar e examinar a consistência técnica do MVP nos cenários definidos.
\section{Métodos}O estudo é aplicado e técnico, baseado em implementação e cenários reproduzíveis. Utilizam-se exclusivamente tarefas sintéticas, sem participantes humanos, contas ou coleta de dados de usuários reais. Cada cenário deve registrar ambiente, versão, comando ou passos, resultado esperado, observado e artefato de evidência. A revisão estática e os testes automatizados examinam regras de tarefa, persistência e backup. A verificação no navegador deve distinguir persistência de dados e carregamento da aplicação offline: salvar tarefas não demonstra que todos os recursos da interface estejam em cache.

Os cenários contemplam criação e conclusão de tarefas, recarga, exportação e importação de backup, rejeição de dados inválidos e reabertura offline após instalação do cache. Resultados são restritos ao ambiente executado. Uma falha deve permanecer registrada; um cenário sem execução não pode ser contado como aprovado.
Flutter oferece execução de aplicações na web; sua documentação informa que versões recentes não geram nem gerenciam um service worker de cache por padrão \cite{flutter}. Portanto, a classificação como PWA e a disponibilidade offline dependem da configuração e da verificação do projeto, não apenas do uso de Flutter.

A propriedade \texttt{localStorage} mantém pares de texto por origem e pode persistir entre sessões \cite{local}. Entretanto, cotas, políticas do navegador, limpeza e modo privado limitam a retenção; armazenamento local não substitui backup \cite{quota}. Um service worker pode interceptar requisições e atender recursos em cache, requer contexto seguro (HTTPS ou localhost) e tem ciclo de instalação e ativação \cite{sw}. Os dados das tarefas e os recursos executáveis precisam ser avaliados separadamente.
\section{Resultados}A execução anterior, v0.1.1, concluiu no GitHub Actions análise estática, seis testes Flutter, quatro testes de versionamento e compilação web. Commit: \texttt{ed3d39a}; execução: \texttt{37249382657}. A reabertura sem rede, a instalação PWA e a falha de quota continuam pendentes de teste específico. Esses resultados técnicos não demonstram eficácia pedagógica \citep{projeto}.

Na demonstração publicada da versão 0.1.0, a inspeção no navegador confirmou criação e conclusão de uma tarefa sintética e recuperação de seu título e estado após recarga online. Uma importação com versão 99 foi rejeitada com mensagem explícita; uma amostra válida restaurou duas tarefas, preservando os estados pendente e concluído. O indicador informou preparação do cache. Isso não comprova recarga sem conexão. A exportação exibiu solicitação de download, mas a captura do arquivo não foi confirmada nesta sessão. Os passos e limites constam de study/browser-observations.md. \citep{projeto}.

\section{Discussão}O artefato constitui pesquisa técnica de MVP, sem revisão por pares, avaliação pedagógica com estudantes, aprovação de banca ou certificação institucional. Não se inferem ganho de aprendizagem, usabilidade populacional ou disponibilidade universal dos testes técnicos. Faltam avaliação em dispositivos variados, sessões privadas, esgotamento de cota e estratégias robustas de atualização. Uma primeira visita sem conexão não deve ser assumida como suportada. A remoção dos dados do navegador pode eliminar tarefas; backup exportado é uma medida de recuperação, cuja integridade também precisa de teste.
\section{Conclusão}O MVP e o protocolo oferecem uma base para revisão técnica escolar. A conclusão sobre cada requisito depende das evidências de execução registradas; a implantação mais ampla requer novas verificações.
\section*{Autoria e uso de IA}Ferramentas de inteligência artificial auxiliaram a elaboração de código e texto nesta demonstração. A atribuição ``Estudantes — identificação omitida nesta demonstração'' não comprova autoria individual. Os estudantes devem conferir fontes, resultados e código, registrar suas contribuições e assumir a versão entregue. Aloisio é indicado como professor/orientador; sua revisão está pendente e não se presume autoria, assinatura ou aprovação.
\begin{thebibliography}{9}
\bibitem{flutter} FLUTTER. \textit{Web FAQ}. Disponível em: \url{https://docs.flutter.dev/platform-integration/web/faq}. Acesso em: 5 out. 2026.
\bibitem{local} MDN. \textit{Window: localStorage property}. Disponível em: \url{https://developer.mozilla.org/en-US/docs/Web/API/Window/localStorage}. Acesso em: 5 out. 2026.
\bibitem{quota} MDN. \textit{Storage quotas and eviction criteria}. Disponível em: \url{https://developer.mozilla.org/en-US/docs/Web/API/Storage_API/Storage_quotas_and_eviction_criteria}. Acesso em: 5 out. 2026.
\bibitem{sw} MDN. \textit{Service Worker API}. Disponível em: \url{https://developer.mozilla.org/en-US/docs/Web/API/Service_Worker_API}. Acesso em: 5 out. 2026.
\bibitem{miller} MILLER, Jane E. Preparing and presenting effective research posters. \textit{Health Services Research}, v. 42, n. 1, p. 311--328, 2007. Disponível em: \url{https://pmc.ncbi.nlm.nih.gov/articles/PMC1955747/}. Acesso em: 5 out. 2026.
\end{thebibliography}
\end{document}
